% from ega4-11-fr.tex
\oldpage[IV-3]{137}

la condition est nécessaire, notons que $\mathcal H$ est de présentation
finie~; la question étant locale sur $X$, on peut supposer $X$ et $Y$ affines,
et, en vertu de (11.3.1), supposer que $\mathcal H$ est $f$-plat. Notons
maintenant le

\begin{lemma}[11.3.9.1]
\label{IV.11.3.9.1-fr}
Soient $f:X\to Y$ un morphisme localement de présentation finie,
$\mathcal G$, $\mathcal H$ deux $\mathcal O_X$-Modules quasi-cohérents de
présentation finie tels que $\mathcal H$ soit $f$-plat,
$p:\mathcal G\to\mathcal H$ un homomorphisme surjectif. Alors $\ker(p)$ est un
$\mathcal O_X$-Module de présentation finie.
\end{lemma}

On peut en effet supposer $X$, $Y$ affines, et qu'il existe un morphisme
$Y\to Y_0$ où $Y_0$ est noethérien, un morphisme $f_0:X_0\to Y_0$ de type fini
tels que $X=X_0\times_{Y_0}Y$, $f=(f_0)_{(Y)}$, deux
$\mathcal O_{X_0}$-Modules cohérents $\mathcal G_0$, $\mathcal H_0$ et un
homomorphisme $p_0:\mathcal G_0\to\mathcal H_0$ tels que $\mathcal G$,
$\mathcal H$ et $p$ se déduisent de $\mathcal G_0$, $\mathcal H_0$ et $p_0$ par
changement de base (8.9.1 et 8.5.2). En outre, on peut supposer $p_0$ surjectif
(8.5.7) et $\mathcal H_0$ $f_0$-plat (11.2.7). Alors, si
$\mathcal K_0=\ker(p_0)$, $\mathcal K_0$ est un
$\mathcal O_{X_0}$-Module cohérent $(\mathbf 0_{\mathrm I},5.3.4)$ et en vertu
de $(\mathbf 0_{\mathrm I},6.1.2)$, $\mathcal K=\ker(p)$ se déduit de
$\mathcal K_0$ par changement de base, donc est de présentation finie.

Ce lemme étant démontré, posons $\mathcal K=\ker(\mathcal G\to\mathcal H)$,
qui est donc de présentation finie. On a un homomorphisme canonique
$q:\mathcal F\to\mathcal K$, et par hypothèse
$q_x:\mathcal F_x\to\mathcal K_x$ est un isomorphisme~; par suite
$(\mathbf 0_{\mathrm I},5.2.7)$ il existe un voisinage $U$ de $x$ tel que
$q|U$ soit un isomorphisme, et en restreignant $X$, on peut supposer que la
suite
\[
0\longrightarrow\mathcal F\xrightarrow{h}\mathcal G
\longrightarrow\mathcal H\longrightarrow0
\]
est \emph{exacte}. Cela étant, il résulte de (11.2.6) que pour un $\lambda$
assez grand, $\mathcal H_\lambda$ est un $\mathcal O_{X_\lambda}$-Module plat~;
posons $\mathcal K_\lambda=\ker(\mathcal G_\lambda\to\mathcal H_\lambda)$ et
$\mathcal K_\mu=v_{\lambda\mu}^*(\mathcal K_\lambda)$ pour $\mu\geq\lambda$~;
il résulte de $(\mathbf 0_{\mathrm I},6.1.2)$ que l'on a
$\mathcal K_\mu=\ker(\mathcal G_\mu\to\mathcal H_\mu)$ et
$\mathcal K=v_\lambda^*(\mathcal K_\lambda)
=\ker(\mathcal G\to\mathcal H)=\mathcal F$. On a d'autre part pour tout
$\mu\geq\lambda$ un homomorphisme canonique
$q_\mu:\mathcal F_\mu\to\mathcal K_\mu$ avec
$q_\mu=v_{\lambda\mu}^*(q_\lambda)$ pour $\lambda\leq\mu$, et
$q=v_\mu^*(q_\mu)$. Comme $q$ est un isomorphisme, il résulte de (8.5.2.4) et
(11.3.9.1) qu'il existe $\mu\geq\lambda$ tel que $q_\mu$ soit un isomorphisme,
et par suite $u_\mu:\mathcal F_\mu\to\mathcal G_\mu$ un homomorphisme
injectif.

Revenons à la démonstration de (11.3.8).

Puisque l'ensemble des $x\in X$ vérifiant b) est ouvert dans $X$, il est clair
que b) entraîne b$'$. Montrons enfin que b$'$ entraîne c). En premier lieu,
$\mathcal F_n$ est $f$-plat dans un voisinage de $x$ (11.3.1) et on peut donc
se borner au cas où $U=X$ et où $\mathcal F_n$ est $f$-plat. Comme par
définition $\operatorname{gr}_{\mathcal J}^*(\mathcal F)$ est isomorphe à
$\mathcal F_n\otimes_{\mathcal O_Y}\mathcal O_Y[T_1,\ldots,T_n]$
$(\mathbf 0,15.2.2)$, il est $f$-plat, et on en conclut par (11.3.4, (i)) que
$\mathcal F$ lui-même est $f$-plat au voisinage de $x$. D'autre part, si
$(\mathcal J_y)_x$ est l'idéal de $\mathcal O_{X_y,x}$ engendré par les
$(g_i\otimes1)_x$, il résulte de (11.2.9.2) que, dans le diagramme
\[
\xymatrix{
(\operatorname{gr}_{\mathcal J_x}^0(\mathcal F_x)[T_1,\ldots,T_n])
 \otimes_{\mathcal O_y}k(y) \ar[r] \ar[d] &
\operatorname{gr}_{\mathcal J_x}^*(\mathcal F_x)
 \otimes_{\mathcal O_y}k(y) \ar[d] \\
(\operatorname{gr}_{(\mathcal J_y)_x}^0((\mathcal F_y)_x))
 [T_1,\ldots,T_n] \ar[r] &
\operatorname{gr}_{(\mathcal J_y)_x}^*((\mathcal F_y)_x)
}
\]

